\chapter{Att kontrollera ett påstående}
\label{app:verifiera}

Det här materialet gör, som allt undervisningsmaterial, en rad
påståenden.  De flesta går att pröva vid datorn: att programmet avbryts
när ett särfall inte tas om hand ser ni själva genom att köra det.  Men
två av påståendena handlar om vad forskningen säger, och dem kan ingen
körning avgöra:
\begin{enumerate}
  \item att studenter systematiskt missförstår felhantering --- att de
    bara använder den allmännaste sorten, och att de inte har tagit till
    sig att resten av \mintinline{python}{try}-blocket hoppas över
    (\cref{sec:fanga,sec:flera});
  \item att det är en känd dålig vana att fånga alla särfall med en enda
    allmän gren (\cref{sec:flera}).
\end{enumerate}
Hur vet man om sådant stämmer?  Samma fråga möter er varje gång ni läser
en lärobok, en webbsida eller ett svar från en språkmodell --- och varje
gång ni själva ska skriva en rapport.  Den här bilagan sammanfattar
arbetsgången; \cref{app:nyborjare,app:catchall} visar den genomförd på
de två påståendena, och \cref{tab:tva-exemplen} sammanfattar frågorna och
svaren.  Arbetsgången beskrivs utförligare i bilagan med samma namn i
föreläsningen \emph{Algoritmiskt tänkande}.

\section{Gör påståendet till en fråga}

Ett påstående går att kontrollera först när det är tydligt vad som skulle
kunna visa att det är \emph{fel}.  Formulera det därför som en fråga med
ett svar.  \enquote{Missförstår studenter felhantering?} kan besvaras med
ja eller nej; \enquote{felhantering är viktigt} kan det inte, eftersom
\enquote{viktigt} är ett omdöme och inte ett faktum.  Ofta döljer sig
flera påståenden i ett: \enquote{en känd dålig vana} är två frågor ---
är den \emph{dålig}, och är den \emph{känd} som sådan --- och de kräver
olika slags källor.

\section{Välj rätt sorts källa}

Olika frågor besvaras av olika källor.
\begin{description}
  \item[Uppslagsverk] för vad ett ord \emph{betyder}.
  \item[Primärkällan] för var en idé \emph{kommer ifrån}, och för ett
    specifikt resultat.  Vill man veta vad studenter faktiskt gjorde
    läser man studien som mätte det, inte en lärobok som återger den.
  \item[Översikter] (\foreignlanguage{english}{reviews}) för om något är
    \emph{etablerat}, eller för vad forskningen \emph{sammantaget}
    säger.
  \item[Enskilda studier] för ett specifikt resultat --- med förbehållet
    att en enskild studie bara visar att något observerats en gång, i en
    viss population.
\end{description}
Var man hittar dem: universitetsbibliotekets söktjänst, och databaser som
Scopus, Web of Science, IEEE Xplore, DBLP (datalogi) och OpenAlex
(öppen).  I bilagorna användes kommandoradsverktyget
\texttt{scholar}\footnote{Paketet \texttt{scholarcli} på PyPI:
  \url{https://pypi.org/project/scholarcli/}.}, som ställer samma fråga
till flera databaser på en gång och sparar allt som hände i en
\emph{session}.  Det ni får ut av bibliotekets söktjänst är detsamma,
bara långsammare.

\section{Sök så att du kan ha fel}

Den som söker på namn hon redan känner till hittar det hon väntar sig.
Sök därför \emph{ämnesdrivet} --- på begrepp, inte författare.  Alla
sökningar i bilagorna är gjorda så.  Några praktiska regler:
\begin{itemize}
  \item Håll frågorna korta.  Flera korta frågor slår en lång.
  \item Ställ \emph{samma} frågor i alla databaser.  Det låter enklare
    än det är: databaserna har olika frågespråk, och en fråga som
    fungerar i en kan ge noll träffar i en annan av rena tekniska skäl.
    \Cref{app:catchall} visar ett exempel --- DBLP indexerar bara
    titlar och kräver att varje ord förekommer, så en lång boolesk
    fråga ger noll där även när databasen \emph{har} rätt artikel.  Det
    är inte ett svar; det är en fråga som ställdes fel, och den ska
    ställas om.
  \item Sök också efter \emph{invändningar}: lägg till ord som
    \enquote{critique}, \enquote{limitations} eller \enquote{justified}.
    Ett påstående som överlevt en motsökning är värt mer än ett som
    aldrig prövats, och en invändning som hittas blir ett
    \emph{förbehåll} i svaret, inte ett nederlag.  \Cref{app:catchall}
    är just ett sådant fall.
  \item Räkna med att databasernas rankning missar somligt.  Vissa
    tidskrifter är inte indexerade alls i de stora databaserna.  Då
    hämtas källan som \emph{känt dokument}: man vet att den finns och
    slår upp den direkt.  Det är i sin ordning, bara det sägs ---
    \cref{app:nyborjare} innehåller ett sådant fall.
\end{itemize}

\section{Läs i fulltext, och skriv ner hur}

En träff i en databas är inte ett belägg.  Sammanfattningen
(\foreignlanguage{english}{abstract}) räcker inte heller: den säger vad
författarna vill ha sagt, inte vad de visat.  Läs hela texten och leta
upp \emph{det ställe} som styrker påståendet --- ett ordagrant citat med
sidnummer.  Hittar ni inget sådant ställe stöder källan inte påståendet,
hur passande titeln än är.  Räkna aldrig en källa ni inte läst.

Skriv sedan ner hur ni gick till väga, så att någon annan kan följa
spåret.  I det här materialet står den anteckningen som en kommentar
ovanför varje post i litteraturlistans källfil, med fasta fält:
\begin{description}
  \item[\texttt{CLAIM}] vilket påstående källan ska styrka;
  \item[\texttt{FOUND-VIA}] hur den hittades --- sökfråga och databas,
    eller \enquote{känt dokument};
  \item[\texttt{PICKED}] varför just den, framför andra träffar;
  \item[\texttt{QUOTE}] det ordagranna citatet, med sida;
  \item[\texttt{VERIFIED}] att fulltexten lästs, och var;
  \item[\texttt{COUNTER}] motsökningen och vad den gav;
  \item[\texttt{DATE}] när.
\end{description}
Det tar ett par minuter per källa, och det är det enda som gör det
möjligt att om ett år svara på frågan \enquote{hur vet du det?}

\section{Dra slutsatsen --- med förbehåll}

Svaret skrivs ut tillsammans med det som talar emot.  Båda bilagorna
svarar \enquote{ja, med förbehåll} (\cref{tab:tva-exemplen}): studenterna
i undersökningarna missförstod mycket riktigt felhanteringen, men de
läste Java och gick sista året, inte Python och första --- så resultatet
säger att svårigheten är verklig och överlever undervisning, inte hur
vanlig den är just hos er.  Och att fånga allt är mycket riktigt en känd
dålig vana, men den studie som säger det säger också att de buggar den
faktiskt orsakar är ovanligare än man kunde tro.  Förbehållen är inte
svagheter i svaret.  De \emph{är} svaret, och de avgör hur påståendet
används i materialet: därför står det i \cref{sec:flera} \enquote{fånga
det ni vet hur ni ska svara på} och inte \enquote{fånga aldrig allt}.
Varje kapitel slutar därför med ett avsnitt \emph{Fortsatt arbete}: vad
just den sökningen lämnade ogjort, och vad som skulle behöva undersökas
för ett säkrare svar.

\begin{table}[htbp]
  \centering
  \caption{De två påståendena som frågor, och svaren i
    \cref{app:nyborjare,app:catchall}.}
  \label{tab:tva-exemplen}
  \begin{tabular}{@{}lp{0.5\linewidth}l@{}}
    \toprule
    Bilaga & Fråga & Svar \\
    \midrule
    \Cref{app:nyborjare} & Missförstår studenter felhantering, och i så
      fall hur? & Ja, med förbehåll \\
    \Cref{app:catchall} & Är det en känd dålig vana att fånga alla
      särfall med en allmän gren? & Ja, med förbehåll \\
    \bottomrule
  \end{tabular}
\end{table}

\section{En anmärkning om språkmodeller}

Sökningarna i bilagorna gav hundratals träffar.  En språkmodell användes
för att \emph{sortera} dem --- hör träffen till ämnet över huvud taget,
och åt vilket håll pekar den? --- så att gallringen kunde granskas.
Lägg märke till arbetsfördelningen: modellen sorterade, men varje källa
som citeras är läst och vald av en människa.  Modellens klassning står
med sin konfidens i tabellerna, så att den kan ifrågasättas.  Använd
gärna språkmodeller för att hitta och sortera; låt dem aldrig vara det
sista ledet.

\chapter{Missförstår studenter felhantering?}
\label{app:nyborjare}
\chapterprecis{Författaren har ännu inte granskat resultaten i den här
  bilagan i sin helhet.}

Materialet påstår på två ställen något om hur studenter tänker.  I
\cref{sec:fanga} sägs att det \enquote{brukar sitta löst} att resten av
\mintinline{python}{try}-blocket hoppas över, och i \cref{sec:flera} att
det är lätt att falla i vanan att bara använda den allmänna sorten.  Båda
är empiriska påståenden om en population, och de styr materialets
utformning: hela kontrasten i \cref{sec:fanga} är byggd kring ett
\mintinline{python}{try}-block med två rader just därför.  Frågan som ska
besvaras är alltså: \emph{missförstår studenter felhantering, och i så
fall hur?}

Att kursens eget underlag om missuppfattningar inte räckte är värt att
säga rakt ut.  Kursens sammanställning över dokumenterade
missuppfattningar hos nybörjare behandlar varken särfall, tracebackar
eller \mintinline{python}{try} --- den närmaste posten gäller
felmeddelanden som en \emph{svårighet}, inte som en missuppfattning.
Påståendena i det här kapitlet är därför sökta från början, inte
återanvända.

\section{Metod}
\label{sec:nyb-metod}

Arbetsgången är den i \cref{app:verifiera}: \emph{finn}, \emph{verifiera},
\emph{registrera}.  Frågan söktes ämnesdrivet och tvåsidigt den 3
september 2026 med verktyget \texttt{scholar}.  Två stödfrågor sökte
belägg för att svårigheterna finns, och två motfrågor sökte arbeten som
säger att de inte finns eller är överdrivna.  Samma begrepp användes åt
båda hållen.  Varje fråga ställdes till alla databaser som var i drift:
OpenAlex, DBLP, IEEE Xplore, Web of Science och Scopus.  Semantic
Scholar avvisade sin nyckel och saknas därför helt.
\Cref{tab:nyb-queries} visar frågorna och antalet träffar per databas.

En källa kunde inte nås med ämnessökning och hämtades som känt dokument.
Rashkovits och Lavys studie av förståelsen av särfallens mer avancerade
egenskaper\autocite{RashkovitsLavy2012} är publicerad i en tidskrift som
ingen av de fem databaserna indexerar; den hittades i stället när
fulltexten till författarnas andra artikel söktes på webben, och lästes
sedan i tidskriftens eget öppna arkiv.  Att ämnessökningen inte gav den
betyder alltså \enquote{inte indexerad här}, inte \enquote{finns inte}.
Databasernas frågespråk krävde också översättning: den booleska strängen
skickades som den var till OpenAlex och IEEE Xplore, Web of Science fick
den som \texttt{TS=(\dots)} och Scopus som
\texttt{TITLE-ABS-KEY(\dots)}, medan DBLP fick en kort ordlista med samma
begrepp.

% Author's note (verification and reproduction): sessions
% exceptions-novice (support S1-S2, counter M1-M2, plus the DBLP keyword
% queries); classification with `scholar llm classify --no-examples`,
% model github_copilot/gpt-5.4; hit-list table generated with
% `scholar sessions export exceptions-novice -f table --bearing-only
% --lang sv --label sok-nyborjare --track "..." -o
% litteratursokning/nyborjare-full`.  The query counts in
% tab:nyb-queries were re-measured per provider with a separate,
% session-free run of the same queries on 2026-09-03.  Bib keys:
% RashkovitsLavy2011, RashkovitsLavy2012, Prather2017CEM.
\begin{table}[p]
  \begin{sidecaption}{Sökfrågor och antal träffar per databas, 3
    september 2026, verktyget \texttt{scholar}.  S = stödfråga,
    M = motfråga, D = samma begrepp som kort ordlista till DBLP.
    Databaser: OpenAlex (OA; inte \foreignlanguage{english}{open
    access}), DBLP, IEEE Xplore (IEEE), Web of Science (WoS), Scopus.
    Högst 40 träffar hämtades per databas och fråga, så 40 betyder
    \enquote{minst 40}.  Frågan visas i den booleska form som skickades
    till OpenAlex och IEEE Xplore.  De långa booleska frågorna ger noll
    i DBLP av tekniska skäl (se \cref{app:verifiera}); D-raderna är
    omsökningen som rättar det.}[tab:nyb-queries]
    \centering\footnotesize
    \begin{tabular}{@{}lp{0.42\linewidth}rrrrr@{}}
      \toprule
      & Nyckelord & OA & DBLP & IEEE & WoS & Scopus \\
      \midrule
      S1 & (exceptions \textsc{or} \enquote{exception handling})
        \textsc{and} (novice \textsc{or} student \textsc{or} CS1
        \textsc{or} \enquote{introductory programming}) \textsc{and}
        (misconception \textsc{or} \enquote{mental model} \textsc{or}
        understanding \textsc{or} difficulties)
        & 40 & 0 & 40 & 40 & 40 \\
      S2 & \enquote{exception handling} \textsc{and} programming
        \textsc{and} (novice \textsc{or} students \textsc{or} learners)
        \textsc{and} (misconception \textsc{or} \enquote{mental model}
        \textsc{or} comprehension)
        & 40 & 0 & 10 & 1 & 2 \\
      M1 & \enquote{exception handling} \textsc{and} (novice \textsc{or}
        students \textsc{or} learners) \textsc{and} (teaching \textsc{or}
        education \textsc{or} curriculum) \textsc{and} (effective
        \textsc{or} successful \textsc{or} \enquote{no difficulty}
        \textsc{or} critique \textsc{or} limitations)
        & 40 & 0 & 40 & 1 & 5 \\
      M2 & \enquote{exception handling} \textsc{and} (students
        \textsc{or} novices \textsc{or} learners) \textsc{and}
        (\enquote{flow of control} \textsc{or} \enquote{control flow}
        \textsc{or} propagation) \textsc{and} (understand \textsc{or}
        mastery \textsc{or} \enquote{not a difficulty} \textsc{or}
        overstated)
        & 40 & 0 & 3 & 0 & 1 \\
      D1 & students exception handling & --- & 1 & --- & --- & --- \\
      D2 & novice exception handling misconceptions
        & --- & 0 & --- & --- & --- \\
      D3 & exception handling education & --- & 0 & --- & --- & --- \\
      \bottomrule
    \end{tabular}
  \end{sidecaption}
\end{table}

\section{Resultat}
\label{sec:nyb-resultat}

Sökningen svarar på frågan i två led.

\paragraph{Använder studenter bara den allmänna sorten?}
Ja.  Rashkovits och Lavy lät 56 studenter lösa en uppgift där
felhantering var poängen, och analyserade sedan vad de faktiskt hade
skrivit.  Resultatet var att \enquote{only few participants (7 out of 56)
provided a solution that was classified to one of the two highest
assimilation levels, while many (23 out of 56) did not use exception
mechanism at all.  The rest of the students (25 out of 56) used the
exception mechanism poorly (i.e., used only Exception class or did not
use hierarchy of exceptions)}\autocite[183]{RashkovitsLavy2011}.  I
intervjuerna beskriver studenterna själva varför: den allmänna sorten är
den enklaste, och meddelandet får skilja fallen
åt\autocite[203]{RashkovitsLavy2011}.

\paragraph{Missförstår de vad som händer med resten av blocket?}
Ja.  I en uppföljande studie ställde samma författare frågor om vad som
händer när ett särfall uppstår mitt i ett block med flera satser.  Om de
studenter som svarade fel skriver de att de \enquote{had not assimilated
the fact that once an exception is thrown, the program jumps to the
closest catch clause that is appropriate for the exception and continues
from there}, och att de därför \enquote{could not cope properly with
try-catch blocks containing more than one
command}\autocite[343]{RashkovitsLavy2012}.  Samma studie räknar upp
\enquote{flow of control in the context of exceptions} och
\enquote{handling exceptions further up the calling chain} bland de
aspekter studenterna hade svårast
för\autocite[350]{RashkovitsLavy2012} --- vilket är precis de två
aspekter \cref{sec:fanga,sec:kasta} är byggda kring.

\paragraph{Är felmeddelandena i sig svåra?}
Ja.  Prather med flera inleder sin studie av hur nybörjare beter sig
inför kompilatorns felmeddelanden med att \enquote{the difficulty in
understanding compiler error messages can be a major impediment to
novice student learning}\autocite{Prather2017CEM}, vilket är det enda
det här materialet stöder sig på från den källan --- motiveringen till
att \cref{sec:felmeddelande} finns alls.

\paragraph{Motsökningen.}
Ingen av de två motfrågorna gav något arbete som säger att studenterna
\emph{inte} har dessa svårigheter, att svårigheten är överdriven eller
att den försvinner med vanlig undervisning.  Frågorna gav träffar
(\cref{tab:nyb-queries}), men de handlade om andra saker: bland de mer
närliggande fanns arbeten om var studenter kör fast i
programvarutestning och om hur AI-assistenter används i
Java-undervisning, ingen av dem med ett motresultat.  Att ingen
invändning hittades är inte samma sak som att ingen finns; det är ett
resultat av en sökning som faktiskt gjordes.

\section{Diskussion}
\label{sec:nyb-diskussion}

Två förbehåll är väsentliga, och båda syns i materialets formuleringar.

Det första är \emph{språket}.  Båda studierna gäller Java.  Java och
Python har samma mekanism --- ett block, en eller flera grenar, ett hopp
till den passande grenen --- men Java tvingar dessutom fram deklarerade
undantag i signaturen, vilket Python inte gör.  Det som mättes är därför
inte identiskt med det som lärs ut här.  Materialet säger detta rakt ut i
\cref{sec:fanga}.

Det andra är \emph{populationen}.  Deltagarna var inte förstaårsstudenter
utan studenter som gått en hel utbildning; den senare studien beskriver
dem som \enquote{on the verge of their graduation}.  Det gör resultatet
starkare i ett avseende och svagare i ett annat.  Starkare, därför att
svårigheten uppenbarligen överlever ordinarie undervisning --- den är
inte något som går över av sig självt.  Svagare, därför att det inte
säger hur vanlig missuppfattningen är i en första kurs.  Vad
materialet därför gör är att \emph{designa} mot svårigheten och låta
övningsfrågorna avslöja hur vanlig den är i just den här gruppen; det är
också vad \cref{ovn:hoppas-over} är till för.

Ett tredje, mindre förbehåll: Prathers studie gäller kompilatorers
felmeddelanden, medan Python är ett tolkat språk.  Materialet drar
slutsatsen i sin svagare form --- att felmeddelanden är svåra --- och
inte i någon starkare.

\section{Slutsats}
\label{sec:nyb-slutsats}

Ja, studenter missförstår felhantering, och de två sätt materialet bygger
på är belagda: de använder den allmänna sorten i stället för att skilja
felen åt, och de har inte tagit till sig att körningen lämnar resten av
blocket när ett särfall uppstår.  Förbehållet är att beläggen kommer från
Java och från studenter i slutet av utbildningen, inte från nybörjare i
Python, och att ingen studie hittades som mäter samma sak i den här
kursens population.  Påståendena bär därför materialets \emph{utformning}
--- vilka kontraster som byggs och vilka frågor som ställs --- men
används inte för att säga åt studenterna vad de tror.

\section{Fortsatt arbete}
\label{sec:nyb-fortsatt}

Sökningen lämnade tre saker ogjorda.  Semantic Scholar avvisade sin
nyckel och kom aldrig med: frågorna ställdes till fem databaser av sex,
och att motsökningen inte gav något väger lättare av just det skälet.
Prathers studie är läst enbart i
sammanfattning\autocite{Prather2017CEM}, och fulltexten återstår att
läsa; den skulle avgöra om resultatet gäller även tolkade språks
felmeddelanden och inte bara kompilatorers --- det förbehåll
\cref{sec:nyb-diskussion} redovisar.  Den tidskrift som publicerade den
senare av de två studierna indexeras inte av någon av de fem
databaserna, utan hämtades som känt dokument; vad som mer står i den
tidskriften vet sökningen alltså ingenting om.  Varje fråga hämtade
dessutom högst 40 träffar per databas, vilket räcker för att se det
högst rankade men inte för att räkna hur mycket som finns.

För ett säkrare svar saknas en studie i den population materialet
vänder sig till.  Ingen av källorna mätte nybörjare i en första kurs,
och ingen mätte Python.  Det som skulle avgöra saken är en upprepning
av Rashkovits och Lavys uppgift och frågor med förstaårsstudenter i
Python: samma uppgift, samma sätt att klassa svaren, en annan
population.  En sådan studie skulle svara på båda de frågor
\cref{sec:nyb-diskussion} lämnar öppna --- om språket spelar roll, och
hur vanlig missuppfattningen är från början.  Visar den att
missuppfattningen är lika vanlig hos nybörjare, kan materialet ange hur
vanlig den är och använda övningsfrågorna som mätning; visar den att
den är ovanlig i Python, krymper kontrasten i \cref{sec:fanga} till ett
exempel och tiden går någon annanstans.

\bigskip
\noindent
\Cref{tab:sok-nyborjare} listar de träffar som rör påståendet.

% --- scholar LaTeX fragment --------------------------------------------
% Generated by: scholar sessions export -f table
% Session: exceptions-novice
% \input this file from the body of a document whose
% preamble loads:
%   \usepackage{longtable}
%   \usepackage{booktabs}
%   \usepackage{array}
% ----------------------------------------------------------------------

\begingroup\footnotesize
\begin{longtable}{@{}%
  >{\raggedright\arraybackslash}p{0.44\textwidth}c%
  >{\raggedright\arraybackslash}p{0.13\textwidth}%
  >{\raggedright\arraybackslash}p{0.26\textwidth}@{}}
\caption{Träffar som rör påståendet, nybörjares förståelse av särfall (2 citerade, 10 som stöder påståendet och 1 som kvalificerar eller motsäger det, av 285 unika träffar; de 134 angränsande och 138 felträffarna listas inte). Skälet till varje inkludering står i sista kolumnen; för maskinklassade rader anges modellens konfidens. Databaser: OpenAlex (OA; inte open access), DBLP, IEEE Xplore (IEEE), Web of Science (WoS), Scopus.}\label{tab:sok-nyborjare}\\
\toprule
Titel & År & Leverantör & Skäl \\
\midrule
\endfirsthead
\multicolumn{4}{@{}l}{\emph{\tablename~\thetable{} (forts.)}}\\
\toprule
Titel & År & Leverantör & Skäl \\
\midrule
\endhead
\midrule \multicolumn{4}{r@{}}{\emph{forts.\ på nästa sida}}\\
\endfoot
\bottomrule
\endlastfoot
Students' misconceptions of java exceptions & 2012 & Scopus & \emph{citerad}: exception-handling \\
Students' Strategies for Exception Handling & 2011 & OA, DBLP, IEEE & \emph{citerad}: exception-handling \\
Coping with abstraction in object orientation with a special focus on application errors & 2010 & OA, DBLP, IEEE & stöder påståendet (0.98) \\
Enhancing Software Testing Education: Understanding Where Students Struggle & 2025 & Scopus & stöder påståendet (0.88) \\
Exception handling negligence due to intra-individual goal conflicts & 2009 & OA, DBLP, IEEE & stöder påståendet (0.96) \\
Learning to Code Together: A Look into Student Projects in an Object-Oriented Programming Course & 2025 & Scopus & stöder påståendet (0.82) \\
Research on the Application of AI Code Assistants in Java Programming Courses & 2026 & OA, DBLP, IEEE & stöder påståendet (0.80) \\
Students' Strategies for Exception Handling & 2011 & WoS & stöder påståendet (0.88) \\
The Development of an Interactive Programming Learning Platform with AI-Driven Feedback and Online Code Execution for the Java Language & 2025 & OA, DBLP, IEEE & stöder påståendet (0.87) \\
Understanding beginners' mistakes with Haskell & 2015 & OA, DBLP, IEEE & stöder påståendet (0.83) \\
Understanding Exception Handling: Viewpoints of Novices and Experts & 2010 & OA, DBLP, IEEE & stöder påståendet (0.98) \\
Use of Java exception stack trace to improve bug fixing skills of intermediate Java learners & 2016 & Scopus & stöder påståendet (0.95) \\
Rethink exception handling at C level & 2010 & OA, DBLP, IEEE & kvalificerar eller motsäger påståendet (0.90) \\
\end{longtable}
\endgroup


\chapter{Är det dåligt att fånga alla särfall?}
\label{app:catchall}
\chapterprecis{Författaren har ännu inte granskat resultaten i den här
  bilagan i sin helhet.}

I \cref{sec:flera} står att ett allmänt \mintinline{python}{except} är en
\enquote{känd dålig vana}.  Det är två påståenden i ett.  Det ena är att
vanan är \emph{dålig} --- att den kostar något.  Det andra är att den är
\emph{känd} som dålig, alltså att det inte är den här kursens privata
åsikt.  Frågan som ska besvaras är: \emph{är det en känd dålig vana att
fånga alla särfall med en allmän gren, och vad kostar den?}

\section{Metod}
\label{sec:ca-metod}

Arbetsgången är den i \cref{app:verifiera}.  Frågan söktes ämnesdrivet
och tvåsidigt den 3 september 2026 med verktyget \texttt{scholar}: två
stödfrågor om vanan som antimönster respektive om tomma hanterare, och
två motfrågor som letade efter arbeten som försvarar vanan eller
begränsar rådet.  Motfrågorna använder samma begrepp som stödfrågorna
plus fältets egna ord för motsatsen (\enquote{justified},
\enquote{intentional}, \enquote{criticism}).  Varje fråga ställdes till
OpenAlex, DBLP, IEEE Xplore, Web of Science och Scopus; Semantic Scholar
avvisade sin nyckel.  \Cref{tab:ca-queries} visar frågorna och antalet
träffar.

Här inträffade det fel \cref{app:verifiera} varnar för.  De fyra
booleska frågorna gav noll träffar i DBLP --- inte därför att DBLP saknar
litteraturen, utan därför att DBLP bara indexerar titlar och kräver att
varje ord förekommer.  Sökningen gjordes därför om mot DBLP med korta
ordlistor med samma begrepp (raderna D1--D3), och först då gav den
träffar, bland dem den artikel om tomma hanterare som stöder
påståendet.  Det är den siffran som står i tabellen.

Klassificeringen av träffarna är delvis maskinell, och en rad har
rättats för hand: priset \enquote{AAPM Coolidge Award, 1976} hade
märkts som stödjande, vilket det uppenbart inte är, och räknas här som
felträff.  Modellens konfidens säger alltså vad modellen trodde, inte
vad som stämmer --- den är ett skäl att läsa raderna, inte ett facit.

% Author's note (verification and reproduction): session
% exceptions-catch-all (support S1-S2, counter M1-M2, known-item K1, DBLP
% redo D1-D3); classification with `scholar llm classify --no-examples`,
% model github_copilot/gpt-5.4; hit-list table generated with
% `scholar sessions export exceptions-catch-all -f table --bearing-only
% --lang sv --label sok-catchall --track "..." -o
% litteratursokning/catchall-full`.  Per-provider counts re-measured in a
% separate session-free run on 2026-09-03.  Bib key: EbertCastor2015.
\begin{table}[p]
  \begin{sidecaption}{Sökfrågor och antal träffar per databas, 3
    september 2026, verktyget \texttt{scholar}.  S = stödfråga,
    M = motfråga, D = samma begrepp som kort ordlista till DBLP.
    Databaser: OpenAlex (OA; inte \foreignlanguage{english}{open
    access}), DBLP, IEEE Xplore (IEEE), Web of Science (WoS), Scopus.
    Högst 40 träffar hämtades per databas och fråga, så 40 betyder
    \enquote{minst 40}.  Frågan visas i den booleska form som skickades
    till OpenAlex och IEEE Xplore; Web of Science fick samma sträng som
    \texttt{TS=(\dots)} och Scopus som
    \texttt{TITLE-ABS-KEY(\dots)}.}[tab:ca-queries]
    \centering\footnotesize
    \begin{tabular}{@{}lp{0.42\linewidth}rrrrr@{}}
      \toprule
      & Nyckelord & OA & DBLP & IEEE & WoS & Scopus \\
      \midrule
      S1 & \enquote{exception handling} \textsc{and} (anti-pattern
        \textsc{or} antipattern \textsc{or} \enquote{code smell})
        \textsc{and} (\enquote{catch all} \textsc{or} \enquote{generic
        exception} \textsc{or} swallow)
        & 40 & 0 & 1 & 0 & 1 \\
      S2 & \enquote{exception handling} \textsc{and} (\enquote{empty
        catch} \textsc{or} \enquote{empty exception handler} \textsc{or}
        \enquote{exception swallowing}) \textsc{and} (practice
        \textsc{or} quality \textsc{or} bug \textsc{or} study)
        & 40 & 0 & 6 & 6 & 7 \\
      M1 & \enquote{exception handling} \textsc{and} (\enquote{catch
        all} \textsc{or} \enquote{generic exception} \textsc{or}
        \enquote{empty catch}) \textsc{and} (justified \textsc{or}
        beneficial \textsc{or} appropriate \textsc{or} intentional
        \textsc{or} \enquote{false positive} \textsc{or} criticism)
        & 40 & 0 & 0 & 0 & 1 \\
      M2 & \enquote{exception handling} \textsc{and} (developers
        \textsc{or} practitioners) \textsc{and} (ignore \textsc{or}
        swallow \textsc{or} suppress) \textsc{and} (rationale \textsc{or}
        deliberate \textsc{or} intentional \textsc{or} perception)
        & 40 & 0 & 0 & 1 & 1 \\
      D1 & empty exception handlers & --- & 1 & --- & --- & --- \\
      D2 & exception handling bugs & --- & 8 & --- & --- & --- \\
      D3 & exception handling empty catch blocks
        & --- & 0 & --- & --- & --- \\
      \bottomrule
    \end{tabular}
  \end{sidecaption}
\end{table}

\section{Resultat}
\label{sec:ca-resultat}

\paragraph{Är vanan känd som dålig?}
Ja, och den har ett etablerat namn.  Ebert, Castor och Serebrenik ---
som undersökt 220 felrapporter och frågat 154 utvecklare --- beskriver i
förbigående, som något självklart, \enquote{overly general catch blocks,
a well-known bad smell in programs that use
exceptions}\autocite[83]{EbertCastor2015}, och hänvisar för namnet
vidare till två äldre arbeten.  Att formuleringen står som en bisats i en
resultatmening är i sig upplysande: det är inte författarnas tes utan
något de förutsätter att läsaren redan vet.

\paragraph{Vad kostar den?}
Mindre än man kunde tro, och det är kapitlets viktigaste fynd.  Samma
mening fortsätter: rapporter om buggar som beror på alltför allmänna
grenar \enquote{are rare, even though there are many opportunities for
them to occur and developers report that they have encountered this kind
of bug in the past}\autocite[83]{EbertCastor2015}.  Studien fann bara 2
av 220 buggar som berodde på tomma hanterare.  Samma studie noterar
också att vanan ofta är medveten: \enquote{it is common for developers to
employ empty catch blocks in circumstances where the implementation of
the system guarantees that the exception will not be
thrown}\autocite[99]{EbertCastor2015}.

\paragraph{Motsökningen.}
De två motfrågorna gav inget arbete som förnekar att vanan är ett
antimönster.  Det som kvalificerar påståendet kom i stället från
stödsidan --- ur samma studie som namnger vanan --- vilket är ett bra
skäl att läsa källor i fulltext i stället för att stanna vid titeln:
titeln lovade stöd, texten gav både stöd och förbehåll.

\section{Diskussion}
\label{sec:ca-diskussion}

Beläggen kommer från stora Java-system (Eclipse och Tomcat) och deras
utvecklare, inte från nybörjarprogram i Python.  Det spelar mindre roll
för \emph{namnet} --- att vanan är känd som dålig är en utsaga om
fältet, inte om ett språk --- men mer för \emph{kostnaden}.  Ett moget
projekt med kodgranskning och statisk analys är inte samma miljö som de
första programmen någon skriver, och \enquote{ovanliga buggar} i Eclipse
säger inte att en nybörjares allmänna gren är ofarlig.  Tvärtom illustrerar
\cref{ex:catchall} en kostnad som studien inte mäter alls: att ett fångat
och omskrivet fel ger användaren \emph{fel} information.

Ett förbehåll om urvalet: de äldre arbeten som Ebert med flera hänvisar
till för namnet har inte lästs i fulltext här.  Påståendet stöder sig
alltså på att en primärstudie behandlar namnet som allmängods, inte på en
egen läsning av dess ursprung.

\section{Slutsats}
\label{sec:ca-slutsats}

Ja, det är en känd dålig vana --- \enquote{a well-known bad smell} är
fältets egen formulering --- men förbehållet är stort nog att ändra
rådet.  De buggar vanan bevisligen orsakar är ovanligare än dess
utbredning antyder, och erfarna utvecklare fångar ibland allt med flit
när de vet att felet inte kan inträffa.  Materialet formulerar därför
regeln som \enquote{fånga det ni vet hur ni ska svara på, och låt resten
synas} och inte som ett förbud, och motiverar den med vad som händer i
\cref{ex:catchall} snarare än med en regel utifrån.

\section{Fortsatt arbete}
\label{sec:ca-fortsatt}

Också den här sökningen saknar Semantic Scholar, som avvisade sin
nyckel: fem databaser av sex svarade.  De två äldre arbeten som Ebert
med flera hänvisar till när de kallar vanan ett känt problem ---
Cabral och Marques från 2007 och Robillard och Murphy från 2003 --- är
varken hämtade eller lästa här, och deras fullständiga uppgifter står i
källförteckningen till den studie kapitlet citerar.  Att läsa dem
skulle visa hur långt tillbaka namnet går och vem som myntade det,
vilket är just det förbehåll \cref{sec:ca-diskussion} redovisar.  De
två citaten på sidorna 83 och 99 återstår dessutom att stämma av mot
den tryckta artikeln.  Omsökningen mot DBLP gjordes med tre korta
ordlistor, och fler formuleringar av samma begrepp skulle rimligen ge
fler träffar där.

Vad som saknas för ett säkrare svar är en mätning i nybörjarnas egen
miljö.  Kostnaden är mätt i stora, mogna Java-system med kodgranskning
och statisk analys, inte i de första programmen någon skriver, och den
kostnad materialet självt visar --- att ett fångat och omskrivet fel
ger användaren fel besked --- mäter studien inte alls.  Det som skulle
avgöra saken är en frekvensstudie i kursens egna inlämningar: hur ofta
en allmän gren förekommer, vad den döljer, och hur lång tid det tar för
studenten att hitta felet ändå.  Visar en sådan studie att den allmänna
grenen sällan kostar något på den här nivån, blir rådet i
\cref{sec:flera} en vana att ta med sig till större program snarare än
en regel för de första; visar den att grenen regelmässigt döljer
skrivfel under felsökning, kan materialet ange kostnaden med källa i
stället för att bara visa den i \cref{ex:catchall}.

\bigskip
\noindent
\Cref{tab:sok-catchall} listar de träffar som rör påståendet.

% --- scholar LaTeX fragment --------------------------------------------
% Generated by: scholar sessions export -f table
% Session: exceptions-catch-all
% \input this file from the body of a document whose
% preamble loads:
%   \usepackage{longtable}
%   \usepackage{booktabs}
%   \usepackage{array}
% ----------------------------------------------------------------------

\begingroup\footnotesize
\begin{longtable}{@{}%
  >{\raggedright\arraybackslash}p{0.44\textwidth}c%
  >{\raggedright\arraybackslash}p{0.13\textwidth}%
  >{\raggedright\arraybackslash}p{0.26\textwidth}@{}}
\caption{Träffar som rör påståendet, att fånga alla särfall (3 citerade, 23 som stöder påståendet och 1 som kvalificerar eller motsäger det, av 133 unika träffar; de 73 angränsande och 33 felträffarna listas inte). Skälet till varje inkludering står i sista kolumnen; för maskinklassade rader anges modellens konfidens. Databaser: OpenAlex (OA; inte open access), DBLP, IEEE Xplore (IEEE), Web of Science (WoS), Scopus.}\label{tab:sok-catchall}\\
\toprule
Titel & År & Leverantör & Skäl \\
\midrule
\endfirsthead
\multicolumn{4}{@{}l}{\emph{\tablename~\thetable{} (forts.)}}\\
\toprule
Titel & År & Leverantör & Skäl \\
\midrule
\endhead
\midrule \multicolumn{4}{r@{}}{\emph{forts.\ på nästa sida}}\\
\endfoot
\bottomrule
\endlastfoot
A Study on Developers\' Perceptions about Exception Handling Bugs. & 2013 & DBLP & \emph{citerad}: handler-generality \\
An exploratory study on exception handling bugs in Java programs. & 2015 & DBLP & \emph{citerad}: citerad \\
Trends on empty exception handlers for Java open source libraries. & 2017 & DBLP & \emph{citerad}: handler-generality \\
A Comparative Study of Error Handling Techniques in Different Programming Languages & 2026 & OA, DBLP, IEEE & stöder påståendet (0.82) \\
A Reflection on \\textquotedbl{}An Exploratory Study on Exception Handling Bugs in Java Programs\\textquotedbl{}. & 2020 & DBLP & stöder påståendet (0.92) \\
AAPM Coolidge Award, 1976 & 1976 & OA, DBLP, IEEE & stöder påståendet (0.90) \\
An Examination of Issues with Exception Handling Mechanisms & 2007 & OA, DBLP, IEEE & stöder påståendet (0.84) \\
Analyzing exception handling in android applications & 2016 & OA, DBLP, IEEE & stöder påståendet (0.92) \\
Characterizing the Exception Handling Code of Android Apps & 2016 & WoS & stöder påståendet (0.98) \\
Designing robust Java programs with exceptions & 2000 & OA, DBLP, IEEE & stöder påståendet (0.90) \\
Discovering faults in idiom-based exception handling & 2005 & OA, DBLP, IEEE & stöder påståendet (0.61) \\
Don't Forget the Exception! : Considering Robustness Changes to Identify Design Problems & 2023 & Scopus & stöder påståendet (0.96) \\
Examining Programmer Practices for Locally Handling Exceptions & 2016 & OA, DBLP, IEEE & stöder påståendet (0.99) \\
Exception handling: A field study in Java and NET & 2007 & WoS & stöder påståendet (0.99) \\
Exception Miner: Multi-language Static Analysis Tool to Identify Exception Handling Anti-Patterns & 2024 & OA, DBLP, IEEE & stöder påståendet (0.95) \\
GLOBAL-AWARE RECOMMENDATIONS FOR REPAIRING EXCEPTION HANDLING VIOLATIONS & 2015 & OA, DBLP, IEEE & stöder påståendet (0.88) \\
How Swift Developers Handle Errors & 2018 & OA, DBLP, IEEE & stöder påståendet (0.88) \\
Illustration and Detection of Exception Handling Bad Smells & 2021 & OA, DBLP, IEEE, Scopus & stöder påståendet (0.96) \\
Revisiting Exception Handling Practices with Exception Flow Analysis & 2017 & OA, DBLP, IEEE & stöder påståendet (0.84) \\
Studying the evolution of exception handling anti-patterns in a long-lived large-scale project & 2020 & OA, DBLP, IEEE & stöder påståendet (0.97) \\
Studying the Prevalence of Exception Handling Anti-Patterns & 2017 & OA, DBLP, IEEE & stöder påståendet (0.99) \\
Supporting Software Quality with Probabilistic Reasoning and Practical Secure Programming Education & 2026 & OA, DBLP, IEEE & stöder påståendet (0.80) \\
The exception handling riddle: An empirical study on the Android API & 2018 & OA, DBLP, IEEE & stöder påståendet (0.86) \\
Toward Declarative Auditing of Java Software for Graceful Exception Handling & 2024 & OA, DBLP, IEEE & stöder påståendet (0.95) \\
Trends on Empty Exception Handlers for Java Open Source Libraries & 2017 & WoS & stöder påståendet (0.99) \\
Unveiling Exception Handling Guidelines Adopted by Java Developers & 2019 & OA, DBLP, IEEE & stöder påståendet (0.93) \\
Challenging Analytical Knowledge On Exception-Handling: An Empirical Study of 32 Java Software Packages & 2014 & OA, DBLP, IEEE & kvalificerar eller motsäger påståendet (0.99) \\
\end{longtable}
\endgroup

